\section{Introduction}
\label{sec:introduction}

A spin cluster carries more information than its total spin. The distribution
of that spin over the sites determines how neighboring clusters interact and
how a momentum-resolved probe couples to their collective motion. In a chain
of clusters, several physical momenta can describe the same effective-chain
momentum, yet their scattering intensities need not be comparable. This makes
spectral weight, rather than the excitation energy alone, a useful probe of
the internal structure of a composite spin.

Trimer chains provide an established example: their ground-state doublets form
an effective spin-$1/2$ chain, while perturbative and numerical spectra connect
higher-energy features to internal trimer excitations~\cite{cheng2022trimer}.
Studies of tetramer chains similarly combine cluster states, effective
descriptions, and dynamical structure factors to distinguish collective and
local excitations~\cite{li2026tetramer}.
The relation between cluster symmetry and the visibility of a transition is
also known for isolated pentamers~\cite{haraldsen2011pentamer}.
These results motivate a specific question for a chain whose cluster spin can
change: how does that change redistribute low-energy intensity among the
physical momenta representing the same collective fluctuation?

We address this question in a linear spin-$1/2$ Heisenberg chain with the
five-bond pattern $J_2$--$J_1$--$J_1$--$J_2$--$J_3$.
The first four bonds define a reflection-symmetric pentamer, and $J_3>0$
couples neighboring cells. We set $J_1=1$ and write
$\alpha=J_2/J_1$ and $\beta=J_3/J_1$.
In the isolated-cell ranges studied, positive $\alpha$ gives a spin-$1/2$
ground doublet, whereas negative $\alpha$ gives a spin-$3/2$ quartet.
Both are half-integer degrees of freedom; the distinction is not the
integer--half-integer contrast between an even and an odd cell.
The Lieb--Schultz--Mattis constraint rules out a unique, symmetry-preserving
gapped bulk ground state for a half-integer spin per translation cell under
its stated symmetry assumptions~\cite{lieb1961,ogata2019}, but it does not fix
the momentum dependence of the response or establish a particular critical
theory throughout this model's coupling plane.

Projecting the five site spins onto the cell ground multiplet gives
$P\bm S_aP=c_a\bm T$. The coefficients $c_a$ determine both the leading
exchange $\Jeff=\beta c_1c_5$ and the interference factor
$F(q)=|\sum_a c_a e^{iqa}|^2$ seen by the probe.
Because the cell momentum is $Q=5q\pmod{2\pi}$, the physical momenta
$q=0.2\pi$, $0.6\pi$, and $\pi$, together with their reflection partners,
all represent $Q=\pi$.
For $\alpha=0.5$ the cell factor strongly favors $q=\pi$, whereas for
$\alpha=-0.5$ it favors $q=0.6\pi$ and nearly cancels at $q=\pi$.
We derive the absolute low-energy spectral weight of two coupled cells and
test it by exact diagonalization over a fixed weak-coupling sequence.
This establishes a quantitative relation between the cell's internal spin
distribution and the relative visibility of folded responses, with an
explicitly tested range of validity.

Reflection provides a second, distinct constraint on the spectrum.
The isolated positive-$\alpha$ cell has an odd-parity lowest excited doublet
that an even $q=\pi$ spin operator cannot reach. Thus the first excited local
level need not be the first visible one.
We compare these local predictions with ground-state and dynamical
matrix-product-state calculations~\cite{white1992,schollwoeck2011,paeckel2019}.
The positive-coupling static data support an effective antiferromagnetic chain
at the tested points. Negative couplings produce a finite-size reconstruction
of the dominant static momentum, and complete connected $L=40$ spectra at
$(\alpha,\beta)=(\pm0.5,0.5)$ compare the two responses at a common
finite-time resolution. These latter points lie outside the tested
weak-coupling regime; they show how spectral weight is organized at finite
coupling, rather than verifying the projected weights there or establishing
a new thermodynamic phase.

Section~\ref{sec:methods} defines the model and numerical observables.
Section~\ref{sec:results} derives and tests the projected spectral weights;
Sec.~\ref{sec:dynamics} presents the finite-coupling chain results.
Section~\ref{sec:discussion} explains the relation between the two local-spin
descriptions and the reconstruction near $\alpha=0$, and identifies the
spectroscopic consequences that follow within the controlled limits.
